\section{Discussion}\label{sec:discussion}

\subsection{What the arms hold constant and what they vary}

The three arms share the same task, the same harness, the same
reflection LM, the same GEPA configuration, and the same convergence
predicate.  What they vary is strictly the evaluator's output
channel:

\begin{itemize}[leftmargin=*,itemsep=2pt,topsep=2pt]
\item \textbf{cert-binary}: $\{0,1\}$ score plus structured
obligation text derived from the theorem's verify function's
evidence dictionary.
\item \textbf{scalar}: graded score on $[0, 1]$ peaking at theorem-holds, no
feedback text beyond the score.
\item \textbf{scalar-evidence}: graded score on $[0, 1]$ peaking at
theorem-holds, plus free-form text enumerating the violating
windows and their means.  No theorem framing.
\end{itemize}

Any observed performance difference, therefore, attaches to one of:
(a) binarization of the score, (b) presence of obligation
evidence text, or (c) presence of the theorem framing line that
names the predicate explicitly.  The active-control arm separates
(a) from (b).  Comparing cert-binary to scalar-evidence isolates
(a) + (c): if the two tie, framing alone did not help; if
cert-binary wins, (a) or (c) carries signal.  The formatter
ablation on seeds 0--4 further separates (b) from
(b) + (c) by stripping the evidence numbers while keeping the
window index, giving three grain levels of text structure.

\subsection{Reading the outcome}

The cert-binary~vs~scalar-evidence comparison is the cleanest
test of the binarization+framing hypothesis, because the two arms
emit byte-identical per-window evidence after the first line ---
only the framing differs (\texttt{Theorem: \dots [FAILED/HOLDS]}
in cert-binary, \texttt{Score: \dots} in scalar-evidence).  In the
reported sweep both arms converge in 1.0 mean iterations across all
ten seeds, Mann-Whitney $U=50$, $p=1.000$, rank-biserial $0.00$ ---
\emph{parity}.  This point estimate is one draw from a
non-deterministic reflection LM; the $N=10$ replication in
Section~\ref{sec:variance} shows the two arms are statistically
indistinguishable in \emph{every} repetition (cert-binary never
slower), so the parity conclusion is robust even though the exact
$p=1.000$ is not.  Under the present LM (Gemma~4, local Ollama) and
task calibration, binarization of the score and explicit naming of
the theorem add no measurable benefit on top of the evidence
content.

The cert-binary~vs~scalar gap, by contrast, is real: in the
reported sweep scalar lags at 2.0 mean iterations ($U=35$,
$p=0.039$, rank-biserial $0.30$), driven by a single seed (seed~6)
that took 9 iterations to converge.  Without any structured failure
information the LM sometimes hunts before finding a stabilising
throttle --- and the replication below shows this hunting is not a
one-seed accident but the scalar arm's defining property: it is the
\emph{high-variance} arm.

The formatter ablation (Table~\ref{tab:paper6-ablation}) refines
this further.  Stripping the per-window numeric evidence down to a
bare index list (``constraint violated: window $k$'') leaves the
default and minimal formatters \emph{tied at 1.0 mean iterations}
on seeds 0--4.  The numeric content of the obligation is not load
bearing on this task --- structural awareness of which window
failed is sufficient for Gemma 4 to choose a throttle low enough to
satisfy the stability predicate in a single step.  Combined with
the cert-binary~vs~scalar-evidence parity, this localises the
signal to a single distinction: \emph{some} structured failure
information vs none.  Neither binarization, theorem framing, nor
numeric evidence content adds further on top.  The replication in
Section~\ref{sec:variance} confirms this ablation tie is not a lucky
draw: default and minimal formatters tie at 1.0 in all ten
warm-LM repetitions.

\subsection{Run-to-run variance}\label{sec:variance}

The reflection LM (\texttt{ollama/gemma4}) is invoked with no
temperature or seed pin, so the per-cell convergence iteration is a
random variable.  The single-sweep point estimates above are
therefore one draw.  To separate a lucky draw from a robust effect
we re-ran the full protocol --- three arms $\times$ ten seeds, plus
the default-vs-minimal ablation on seeds 0--4 --- $N=10$ times
(\texttt{eval/convergence/theorem\_6\_variance.py}, aggregate at
\texttt{eval/results/theorem\_6/variance\_summary.json}).  Each
repetition holds GEPA's \texttt{seed} and the SHA-256 data seeding
fixed, so the only varying factor is LM sampling.  Reported as
median~[min,~max] of the per-repetition arm mean:

\begin{itemize}[leftmargin=*,itemsep=2pt,topsep=2pt]
\item \textbf{cert-binary}: $1.0~[1.0,~1.0]$ --- zero variance; every
seed converges in one step in every repetition.
\item \textbf{scalar-evidence}: $1.0~[1.0,~1.1]$ --- near-zero
variance; a single seed flipped to two iterations in one of ten
repetitions.
\item \textbf{scalar}: $1.7~[1.2,~6.5]$ --- the high-variance arm.
One repetition recorded a non-converged run (a seed exhausting the
50-iteration budget); the graded-reward-only signal leaves the LM
hunting unpredictably.
\end{itemize}

Two conclusions are \emph{robust} across all ten repetitions, and
two point estimates are \emph{not}.  Robust: (i) cert-binary and
scalar-evidence are statistically indistinguishable --- Mann-Whitney
$p \ge 0.05$ in $10/10$ repetitions (observed $p$ range
$[0.18,~1.00]$), with cert-binary never slower (ratio $\le 1.0$
throughout); and (ii) the predeclared $\le 0.75$-ratio gate is
cleared in $0/10$ repetitions, so the Null verdict on Theorem~3 does
not depend on the reported sweep.  Not robust: the exact $p=1.000$
and the exact scalar mean of $2.0$ are single-draw artifacts.

The variance itself carries the most interesting signal.  Structured
failure information does not merely lower the mean iteration count;
it \emph{collapses the variance}.  Both cert-binary and
scalar-evidence are essentially deterministic at 1.0, while the
graded-scalar arm ranges over a factor of five and occasionally
fails to converge at all.  On this task the value of an obligation
channel --- certificate or evidence text --- is better described as
\emph{reducing the variance of convergence} than as reducing its
mean.  This reframing is consistent with the Null verdict: certificates
buy reliability that a graded scalar with the same evidence text
already supplies, not a speed-up over it.

\subsection{What this does not show}

The synthetic task is a \emph{signal-channel} test, not a
general-capability test.  A cert-binary win on a one-knob task does
not establish that certificate-based evaluators dominate on
multi-component DSPy programs, on non-enumerable failure modes, or
on tasks where the reflection LM is weaker at parsing structured
constraints than numeric gradients.

Our LM is Gemma 4, run locally via Ollama.  Frontier models may
behave differently, particularly on the scalar-evidence arm where
the feedback is rich English.  A cross-LM robustness follow-up is
declared as out of scope here; see Section~\ref{sec:followups}.

\subsection{Relationship to Operon's broader thesis}\label{sec:thesis}

The Operon convergence layer already supports six external
frameworks at the \emph{topology} level (swarms, DeerFlow,
AnimaWorks, Ralph, A-Evolve, Scion).  This paper's contribution is
to extend convergence to the \emph{evaluator} level: a single
integration module (\texttt{operon\_ai/convergence/gepa\_adapter.py})
lets any of Operon's eight registered theorems stand in for
GEPA's scalar reward.  Whether that integration \emph{helps} is the
empirical question this paper answers.  The theorem
registry~\cite{banu2026operon} is append-only: downstream users register a
verify function once via
\texttt{operon\_ai.core.certificate.register\_verify\_fn} and the
GEPA adapter picks it up automatically.
